% \section{Related Work}
% \label{sec:related_work}

% This section reviews prior work across three key areas relevant to our research: (1) automated data preparation, (2) data synthesis methods for training data preparation models, and (3) training methodologies for complex reasoning tasks. We situate \sys within this landscape, highlighting its unique contributions to end-to-end, NL-\nlwhat data preparation.

% \subsection{Data Preparation}

% The field of data preparation has seen significant research, largely branching into example-driven systems and task-specific solutions.

% \stitle{Example-Driven Transformation.} A prominent line of work leverages programming-by-example (PBE), where users provide one or more input-output examples to specify the desired transformation~\cite{barowy2015flashrelate, singh2016blinkfill, he2018transform, gulwani2012spreadsheet, jin2017foofah, heer2015predictive}. For instance, FLASHRELATE~\cite{barowy2015flashrelate} infers data extraction rules from user-provided examples by considering cell content and spatial constraints. However, these methods shift the burden to the user to provide representative examples. In the context of end-to-end data preparation, where the final target table is the result of a long and complex operation chain, constructing a valid and comprehensive set of examples is often impractical or as difficult as writing the program itself.

% \stitle{Sub-Task Specific Solutions.} The majority of existing automated data preparation systems focus on solving isolated sub-problems rather than the end-to-end challenge. These efforts typically fall into one of three categories: Data Cleaning~\cite{chen2025auto, huang2018auto, fan2024cost}, Column Transformation~\cite{he2018transform,jin2017foofah,fan2024autoprep}, or Table Transformation~\cite{zhu2017auto, DBLP:journals/pvldb/LiHYWC23}. For example, Auto-Tables~\cite{DBLP:journals/pvldb/LiHYWC23} and Auto-Prep~\cite{lai2025auto} focus specifically on restructuring non-relational tables through operations like Pivot, Transpose, and Unpivot. While these tools achieve high performance on their designated tasks, they are not designed to compose operations across different categories and thus fail to address the holistic, multi-step nature of real-world data preparation workflows that require seamless orchestration of cleaning, transformation, and restructuring.

% The most relevant work to ours is Text2Pipeline~\cite{ge2025text}, which established a benchmark for controlling data preparation processes using natural language. However, our approach diverges from and improves upon Text2Pipeline in two fundamental ways. \etitle{(1) Input Modality.} Text2Pipeline requires the user to provide a comprehensive and precise natural language description of the entire transformation process. In practice, articulating such a detailed sequence of steps is challenging and unnatural for users. \sys adopts a more flexible and intuitive approach where the user primarily provides a description of the target table schema. The system then infers the necessary intermediate operations. This aligns better with how users think—in terms of the desired outcome, not the procedural steps. \etitle{(2) Operator Flexibility.} Our operator space is more comprehensive and robust. Crucially, it supports passing functions as parameters (\eg a lambda function for a complex SplitColumn operation), offering a degree of flexibility that is absent in systems with fixed parameter sets.

% \subsection{Data Synthesis for Data Preparation}

% A major bottleneck in training robust data preparation models is the scarcity of full-chain supervision data. Existing data synthesis methods are often ill-suited for the complexities of data preparation.

% \stitle{Task-Driven Synthesis.} In related domains like NL2SQL, task-driven synthesis has been successful. For example, OmniSQL~\cite{li2025omnisql} generates SQL queries of varying complexity from a database schema and then uses back-translation to create corresponding natural language questions. While effective for generating question-query pairs, this approach is less suitable for data preparation. DP tasks involve long, procedural chains where the validity of an operation depends on the state of the data from the previous step. Simply generating a question and a final solution (\eg a target table) fails to produce the critical intermediate operation sequence needed for supervision.

% \stitle{Rule-Based Inverse Synthesis.} A more related approach is found in Auto-Tables~\cite{DBLP:journals/pvldb/LiHYWC23}, which synthesizes data by applying a sequence of predefined inverse operations to a clean relational table to generate a non-relational source table. For example, the inverse of wide-to-long is defined as a sequence of stack, split, and pivot. While innovative, this method has two key limitations. First, it is restricted to operations with fixed parameters and easily definable inverses (\eg transpose is its own inverse), making it inapplicable to more complex, data-dependent operations in our space. Second, its reliance on fixed rules introduces a strong distributional bias into the training data, failing to capture the diversity and unpredictability of data quality issues found in the real world. In contrast, our bi-directional synthesis pipeline leverages LLMs for both forward (SQL-to-operations) and backward (clean-to-dirty) generation, enabling the creation of diverse, realistic, and verifiable full-chain training examples.

% \subsection{Training Methods for Data Preparation}

% The training paradigm is critical for enabling a model to handle complex, multi-step reasoning.

% \stitle{Supervised Fine-Tuning.} Current training paradigms for data preparation models, such as JellyFish and TableGPT, are predominantly based on Supervised Fine-Tuning (SFT). SFT is effective for teaching models to memorize mappings for single-step operations or short, common sequences. However, it falls short in developing the advanced reasoning, planning, and error-correction capabilities required for long-chain DP tasks. The exponential search space of possible operation sequences means that SFT alone cannot equip a model to robustly explore, backtrack, and find correct solutions for novel problems.

% \stitle{Reinforcement Learning for Reasoning.} To overcome these limitations, we turn to Reinforcement Learning (RL), which has proven effective in enhancing the reasoning and planning abilities of LLMs~\cite{RL work}. To the best of our knowledge, \sys is the first framework to apply multi-turn RL to the general, end-to-end data preparation task. Unlike prior work that may use single-turn RL, our multi-stage policy learning framework, which combines a progressive cold start with multi-turn RL, is specifically designed to manage long-chain orchestration. By introducing a hybrid reward function that provides dense, partial credit and semantic feedback, we enable the model to effectively learn complex exploration and rollback strategies, navigating the vast search space to achieve state-of-the-art performance.

\section{Related Work}
\label{sec:related_work}

\stitle{Data Preparation.}
Prior data preparation work mainly focuses on programming-by-example (PBE).
PBE systems~\cite{barowy2015flashrelate, singh2016blinkfill, he2018transform, gulwani2012spreadsheet, jin2017foofah} synthesize transformation programs from input-output examples, which are effective for well-specified transformations but require example-level supervision for each task.
Other systems target specific sub-problems of data preparation ~\cite{yang2021auto, fan2024autoprep, lai2025auto, chen2025auto, DBLP:journals/pvldb/LiHYWC23,huang2018auto, fan2024cost}. For example, Auto-Tables~\cite{DBLP:journals/pvldb/LiHYWC23} focuses on table reconstruction operations such as pivoting and stacking. Such systems 
emphasize structural transformations but
are not designed to jointly handle value-level cleaning~\cite{chai2025cost} and multi-table integration~\cite{zhu2017auto}. Text-to-Pipeline~\cite{ge2025text} maps natural language descriptions to pipelines, but assumes detailed step-wise specifications as input.
In contrast, \sys targets end-to-end autonomous data preparation from high-level target schema descriptions and supports a much broader operator space covering cleaning, reshaping, and integration.

\stitle{LLM-Powered Data Preparation.}
Recent work explores using LLMs for data preparation pipeline construction. Prompting-based approaches directly generate transformation code from natural language specifications~\cite{li2024towards,sapkota2025vibe,lai2025auto}. 
To incorporate execution feedback, ReAct-style methods interleave reasoning with tool actions and update subsequent decisions based on observed results~\cite{yao2023react, aksitov2023rest}. Such approaches follow a linear interaction trajectory, where actions are appended sequentially. More recent work explores tree-structured search over reasoning or operator spaces, including Tree-of-Thoughts (ToT)~\cite{yao2023tree} and MCTS-based agents~\cite{li2025alphasql, ge2025monteprep}. For example, MCTS-style approaches typically rely on rollout-based value estimates and scalar reward signals to guide node selection and backpropagation~\cite{hao2023reasoning}. However, data preparation pipelines expose structured execution feedback, such as schema mismatches, execution errors for operators, and intermediate table states. \sys is designed to operate directly over such execution-grounded states, where each node corresponds to materialized tables and branch decisions are guided by structured execution feedback rather than scalar rollout statistics. This design supports operator-level revision and state-aware backtracking during pipeline construction.

% \stitle{Data Synthesis for Data Preparation.}
% Training data for preparation is limited, especially for tasks requiring multi-step data preparation pipelines. NL2SQL synthesis methods~\cite{li2025omnisql} generate query–result pairs, but typically do not provide operator-level transformation sequences needed for ADP training. Auto-Tables~\cite{DBLP:journals/pvldb/LiHYWC23} proposes a rule-based inverse-operator synthesis strategy for reshaping tasks, mainly targeting structural transformations, which can limit the diversity of synthesized data. Our synthesis framework differs in scope by generating operator pipelines together with source–target table pairs and reversible noise injection, supporting training signals for autonomous data preparation.



%\stitle{Automated Data Preparation.}
%Prior work has largely focused on either programming-by-example (PBE) or isolated sub-tasks. PBE systems~\cite{barowy2015flashrelate, singh2016blinkfill, he2018transform, gulwani2012spreadsheet} require users to provide input-output examples, which is impractical for specifying complex, multi-step transformations. 
%Concurrently, many systems address specific sub-problems. For instance, Auto-Tables~\cite{DBLP:journals/pvldb/LiHYWC23} establishes state-of-the-art performance in \textit{table reshaping} (\eg transforming non-relational headers via Pivot/Stack). However, it is limited to structural alignment and lack the capability to handle value-level anomalies (\eg missing values, typos) or multi-table integration (\eg joins). 
%\sys advances this by supporting a comprehensive library of 31 operators covering cleaning, reshaping, and integration. Unlike Text2Pipeline~\cite{ge2025text} which requires detailed procedural descriptions, \sys infers the entire pipeline solely from the target schema, significantly reducing user burden.

%\stitle{Data Synthesis for Data Preparation.}
%Training robust ADP models is hindered by the scarcity of full-chain supervision data. 
%While NL2SQL synthesis~\cite{li2025omnisql} generates query pairs, it fails to produce the intermediate operation sequences crucial for ADP. 
%Auto-Tables~\cite{DBLP:journals/pvldb/LiHYWC23} pioneered a rule-based inverse synthesis approach, generating training data by applying the mathematical inverse of reshaping operators (\eg applying Unpivot to simulate a source table). 
%While effective for structural layout tasks, this rule-based method cannot simulate realistic, semantic data quality issues (\eg inconsistent date formats or dirty categorical values) required for data cleaning tasks.
%Our bi-directional synthesis overcomes this limitation by leveraging LLMs to generate context-aware inverse logic. This allows \sys to inject diverse, recoverable noise into clean datasets, covering both complex structural changes and realistic data dirtiness.

%\stitle{Agentic Reasoning and Search Strategies.}
%LLMs have evolved from simple code generators to autonomous agents. 
%Early agents adopted linear reasoning frameworks like ReAct~\cite{yao2023react} or Chain-of-Thought, where actions are appended sequentially. 
%However, data preparation is inherently interdependent where early mistakes often manifest as errors only after several steps. Linear agents struggle to recover from such non-local errors without restarting.
%To address this, recent works introduce Tree-based Search methods, such as Tree of Thoughts (ToT)~\cite{yao2024tree} and Monte Carlo Tree Search (MCTS)~\cite{li2025alphasql, ge2025monteprep}. 
%While MCTS enables backtracking, it is suboptimal for data preparation for two reasons: 
%(1) Local Blindness: MCTS expands nodes based on local state statistics, ignoring valuable failure signals from sibling branches (\eg a specific cleaning step failing in a parallel branch); 
%(2) Scalar Feedback: It relies on scalar rewards (\eg pass/fail rates) for backpropagation, discarding the rich semantic feedback (\eg specific error traces or data anomalies) provided by the execution environment.
%In contrast, \sys employs a \textit{Agentic Reasoning Tree} that explicitly maintains the history of all branches and utilizes execution feedback to guide precise, semantic-aware backtracking and error correction.



% 重点和Auto-Table做比较